% Source PDF: source/普通高考/2015/2015安徽理.pdf, page 1, question 13.
% The flowchart is vector line art; pdfimages -list found no embedded bitmap objects.
% Grid evidence: tmp/review_flowchart_2015_anhui_science/grid/page-1-grid100.png
% (100px coarse + 50px fine) and q13-block-grid50.png (50px + 25px fine).
% Original comparison crop: tmp/review_flowchart_2015_anhui_science/crop/q13_original_final.png,
% cropped from the ungridded page render.
% Elements: rounded start/end terminals, initialization and update rectangles, a decision diamond,
% output trapezoid, directed connectors, branch labels, and the left return loop.
% Node -> directed edge list: 开始→(a=1,n=1)→判定；判定(是)→(a=1+1/(1+a))→(n=n+1)→回到判定；
% 判定(否)→输出n→结束。回边经判定上方单一汇合点返回，不与输出支路共享或交汇。
% solid segments: all node borders and directed connectors; dashed segments: none.
% labels: formulas are centered in nodes; 是 is beside the downward branch and 否 above the right branch.
\begin{tikzpicture}[
  x=1cm,y=1cm,
  >=Stealth,
  line width=0.55pt,
  line cap=round,line join=round,
  every node/.style={font=\normalsize,anchor=center,align=center,
    execute at begin node={\setlength{\parindent}{0pt}}},
  flowbox/.style={draw,minimum height=.68cm,inner xsep=4pt,inner ysep=2pt,
    anchor=center,align=center,
    execute at begin node={\setlength{\parindent}{0pt}}},
  io/.style={draw,shape=trapezium,trapezium left angle=78,trapezium right angle=102,
    trapezium stretches=false,minimum height=.78cm,inner xsep=4pt,inner ysep=2pt,
    anchor=center,align=center,
    execute at begin node={\setlength{\parindent}{0pt}}},
  flowarrow/.style={->,line width=0.55pt},
]
  \node[draw,rounded corners=7pt,minimum width=1.35cm,minimum height=.72cm]
    (start) at (0,8.55) {开始};
  \node[flowbox,minimum width=2.30cm] (init) at (0,7.18) {$a=1,n=1$};
  % Keep the merge comfortably below the initialization box (>=2.5 mm).
  \coordinate (merge) at (0,6.50);
  \node[draw,diamond,aspect=2.85,minimum width=5.55cm,minimum height=1.12cm,
    inner sep=2pt,anchor=center,align=center,
    execute at begin node={\setlength{\parindent}{0pt}}] (test) at (0,5.62)
    {$\lvert a-1.414\rvert\geq 0.005?$};
  \node[flowbox,minimum width=2.45cm,minimum height=.92cm] (update) at (0,4.02)
    {$a=1+\dfrac{1}{1+a}$};
  \node[flowbox,minimum width=1.65cm] (inc) at (0,2.58) {$n=n+1$};
  \node[io,minimum width=1.80cm] (output) at (4.50,4.02) {输出 $n$};
  \node[draw,rounded corners=7pt,minimum width=1.05cm,minimum height=.72cm]
    (finish) at (4.50,2.68) {结束};

  % Main path and the single merge point above the decision.
  \draw[flowarrow] (start.south) -- (init.north);
  \draw (init.south) -- (merge);
  \draw[flowarrow] (merge) -- (test.north);

  % 是 branch, update, increment, and independent loop return to the merge point.
  \draw[flowarrow] (test.south) -- (update.north);
  \draw[flowarrow] (update.south) -- (inc.north);
  \coordinate (loop-bottom) at (0,1.38);
  \coordinate (loop-left) at (-3.05,1.38);
  \coordinate (loop-entry) at (-3.05,6.50);
  \draw (inc.south) -- (loop-bottom) -- (loop-left) -- (loop-entry);
  % Keep the loop arrowhead clear of the degree-3 merge junction.
  \draw[flowarrow] (loop-entry) -- (-0.35,6.50);
  \draw (-0.35,6.50) -- (merge);

  % 否 branch to output and finish.
  \draw[flowarrow] (test.east) -- (4.50,5.62) -- (4.50,4.42);
  \draw[flowarrow] (output.south) -- (finish.north);

  \node[inner sep=0pt,anchor=west] at (0.65,4.78) {是};
  \node[inner sep=0pt,anchor=west] at (3.45,6.12) {否};

  \path[use as bounding box] (-3.60,8.95) rectangle (5.55,0.95);
\end{tikzpicture}
